%--- iliad ---
% title: Training Dynamics
% summary: >-
%   Exact learning dynamics of deep linear networks -- loss-landscape geometry,
%   balanced gradient flow and the NTK, the rich (saddle-to-saddle) and lazy
%   regimes, their mixed unification, and the implicit bias of SGD noise.
%--- end ---
\documentclass[11pt]{article}

% iliad.sty: local copy first (Overleaf), else the shared ../iliad.sty.
\IfFileExists{iliad.sty}{\usepackage[boxes]{iliad}}{\usepackage[boxes]{../iliad}}

% ---- Free zone -------------------------------------------------------------
% tightlist is emitted by the pandoc source lists we copied verbatim.
\providecommand{\tightlist}{%
  \setlength{\itemsep}{0pt}\setlength{\parskip}{0pt}}
% Unicode glyphs that appear verbatim in the copied prose/table.
\usepackage{newunicodechar}
\newunicodechar{→}{\ensuremath{\rightarrow}}

\title{Training Dynamics}
\author{\authorname{Guillaume Corlouer}\\ \affiliation{Stormglass}}
\date{\today}

\begin{document}
\maketitle

% Guillaume's recorded lectures for this day (titles auto-queried at build).
\youtube{IHHrJdZKICs}   % deep linear networks
\youtube{clTLzFD2uC8}   % learning dynamics

\begin{learningoutcomes}
  \subsection*{Motivation}
  \begin{itemize}
    \item Know about some big open questions in learning dynamics
    \item Understand the concept of implicit regularization
    \item Know about the different approaches to studying learning dynamics
    \item Understand the AI safety motivations for learning dynamics
    \item Know about emergent misalignment as a safety-relevant phenomenon that illustrates the importance of understanding generalization for AI safety
  \end{itemize}

  \subsection*{Geometry and dynamics of deep neural networks}
  \begin{itemize}
    \item Know key results about the loss landscape of deep linear networks (DLNs): critical points are saddles or global minima
    \item Explain the edge of stability phenomenon
    \item Understand that gradient flow can be written as NTK-weighted gradient in function space
    \item DLNs are degenerate and have conserved quantities through gradient flow
  \end{itemize}

  \subsection*{Regimes of learning in deep linear networks}
  \begin{itemize}
    \item Understand the role of initialization, width and depth for the lazy and rich (saddle to saddle) regimes in DLNs
    \item Know about implicit regularization from non-linearity (neural race reduction)
  \end{itemize}
\end{learningoutcomes}

Use LLMs to help you out when you spend longer than the suggested time (but make sure that you understand). Make sure you keep at least 30 minutes to do problem 3.

\textbf{References}: \cite{saxe2014}, \cite{achour2024}, \cite{tu2024}.

\section{Setup and notation}

We study deep linear networks (DLNs) with \(L\) weight matrices
\(W_1 \in \mathbb{R}^{d_1 \times d_0}, W_2 \in \mathbb{R}^{d_2 \times d_1}, \ldots, W_L \in \mathbb{R}^{d_L \times d_{L-1}}\).
The network computes:

\[f(x) = W_L W_{L-1} \cdots W_1 \, x =: W x\]

where \(W = W_L \cdots W_1 \in \mathbb{R}^{d_L \times d_0}\) is the
end-to-end (or ``student'') matrix. We train on a dataset
\(\{(x_\mu, y_\mu)\}_{\mu=1}^N\) with the squared loss:

\[\mathcal{L}(\theta) = \frac{1}{2N} \sum_{\mu=1}^{N} \|y_\mu - f(x_\mu)\|^2\]

In the population limit \(N \to \infty\) with whitened inputs \(\Sigma_X = I\), this becomes, up to an additive constant\footnote{The constant is \(\frac{1}{2}\mathbb{E}\|y - Mx\|^2\), which is the residual of the OLS. It vanishes on realizable data.}:

\[\mathcal{L}(W) = \frac{1}{2}\|M - W\|_F^2\]

where \(M = \Sigma_{YX} \Sigma_X^{-1} = \Sigma_{YX}\) is the ``teacher'' matrix (the
OLS solution) with SVD
\(M = U \,\text{diag}(s_1, \ldots, s_r, 0, \ldots, 0)\, V^\top\), and
\(s_1 \geq s_2 \geq \cdots \geq s_r > 0\).

% ---------------------------------------------------------------------------
% %%% Where the additive constant comes from.
%
% In the population limit the loss is \frac{1}{2}\mathbb{E}\|y - Wx\|^2, with
% second moments \Sigma_X = \mathbb{E}[x x^\top] and
% \Sigma_{YX} = \mathbb{E}[y x^\top].
%
% Split the residual into
%
%     y - Wx = (y - Mx) + (M - W)x.
%
% Note \mathbb{E}[(y - Mx)x^\top] = \Sigma_{YX} - M\Sigma_X = 0.
% The cross term obtained from expanding the square in the loss
% is \text{Tr}[(M - W)\,\mathbb{E}[x(y - Mx)^\top]] and so
% vanishes, leaving
%
%     \mathcal{L}(W) = \frac{1}{2}\mathbb{E}\|(M - W)x\|^2 + \mathcal{L}_\star,
%     \qquad \mathcal{L}_\star := \frac{1}{2}\mathbb{E}\|y - Mx\|^2.
%
% Whitened inputs, \Sigma_X = I, then reduce the first term to the Frobenius
% norm, giving \mathcal{L}(W) = \frac{1}{2}\|M - W\|_F^2 + \mathcal{L}_\star.
% ---------------------------------------------------------------------------

\section{Loss landscape geometry}
\label{sec:td-geometry}

This problem explores the critical point structure of deep linear
networks, following \cite{achour2024}.

\begin{exercise}[Diagonal decomposition]
\label{ex:td-1a}
Consider the \textbf{diagonal} case:
\(d_0 = d_1 = \cdots = d_L = d\), and the teacher is diagonal,
\(M = \text{diag}(s_1, \ldots, s_d)\), with
\(s_1 > s_2 > \cdots > s_d > 0\). Restrict attention to diagonal weight
matrices \(W_l = \text{diag}(w_l^{(1)}, \ldots, w_l^{(d)})\). Show that
the loss decomposes into \(d\) independent scalar problems:

\[\mathcal{L} = \frac{1}{2}\sum_{\alpha=1}^d \left(s_\alpha - \prod_{l=1}^L w_l^{(\alpha)}\right)^2\]
\end{exercise}

\begin{solution}[ex:td-1a]
When all \(W_l\) are diagonal, the product
\(W = W_L \cdots W_1\) is also diagonal with entries
\((W)_{\alpha\alpha} = \prod_{l=1}^L w_l^{(\alpha)}\). The teacher \(M\)
is diagonal with entries \(s_\alpha\). Then:

\[\mathcal{L} = \frac{1}{2}\|M - W\|_F^2 = \frac{1}{2}\sum_{\alpha=1}^d (s_\alpha - (W)_{\alpha\alpha})^2 = \frac{1}{2}\sum_{\alpha=1}^d \left(s_\alpha - \prod_{l=1}^L w_l^{(\alpha)}\right)^2\]

Since the different modes \(\alpha\) share no parameters, the loss
decomposes into \(d\) independent scalar problems. \(\square\)
\end{solution}

\begin{exercise}[Scalar critical points]
\label{ex:td-1b}
For a single scalar mode with target \(s > 0\), find all
first-order critical points of
\(\ell(w_1, w_2) = \frac{1}{2}(s - w_1 w_2)^2\) at depth \(L = 2\). Show
that the critical points are:

\begin{itemize}
\tightlist
\item
  The global minimum manifold: \(w_1 w_2 = s\)
\item
  The origin: \(w_1 = w_2 = 0\)
\end{itemize}

Classify the origin as a saddle point by computing the Hessian \(H\) of
\(\ell\) at \((0,0)\) and showing it has both positive and negative
eigenvalues.

\emph{Hint}: Write the gradient equations.
\end{exercise}

\begin{solution}[ex:td-1b]
For \(L = 2\), the gradient equations are:

\[\frac{\partial \ell}{\partial w_1} = -(s - w_1 w_2)\,w_2 = 0, \qquad \frac{\partial \ell}{\partial w_2} = -(s - w_1 w_2)\,w_1 = 0\]

\textbf{Case 1}: \(s - w_1 w_2 = 0\), i.e.~\(w_1 w_2 = s\). This is the
global minimum manifold (a hyperbola \(w_2 = s/w_1\)) with \(\ell = 0\).

\textbf{Case 2}: If \(s - w_1 w_2 \neq 0\), then the first equation
requires \(w_2 = 0\) and the second requires \(w_1 = 0\). (For instance,
if \(w_2 \neq 0\) and \(w_1 = 0\), the first equation gives
\(-s\,w_2 = 0\), which contradicts \(s > 0\) and \(w_2 \neq 0\).) So the
only other critical point is \(w_1 = w_2 = 0\).

The Hessian at \((0,0)\). We need the second derivatives of
\(\ell = \frac{1}{2}(s - w_1 w_2)^2\):

\begin{align*}
\frac{\partial^2 \ell}{\partial w_1^2} &= w_2^2, \qquad \frac{\partial^2 \ell}{\partial w_2^2} = w_1^2, \\
\frac{\partial^2 \ell}{\partial w_1 \partial w_2} &= -(s - w_1 w_2) + w_1 w_2 = 2w_1 w_2 - s
\end{align*}

At \((0,0)\):

\[H = \begin{pmatrix} 0 & -s \\ -s & 0 \end{pmatrix}\]

The eigenvalues are \(\pm s\). Since \(s > 0\), \(H\) has one positive
and one negative eigenvalue. The origin is a \textbf{strict saddle
point}. \(\square\)
\end{solution}

\begin{exercise}[Critical-point structure]
\label{ex:td-1c}
Achour et al.~\cite{achour2024} show that every
first-order critical point \(\theta = (W_1, \ldots, W_L)\) satisfies:
there exists a subset \(S \subseteq \{1, \ldots, d\}\) such that

\[W = W_L \cdots W_1 = P_S\, M\]

where \(P_S = U_S U_S^\top\) is the orthogonal projector onto the span
of the left singular vectors \(\{u_\alpha\}_{\alpha \in S}\) of \(M\).
For the diagonal case \(M = \text{diag}(s_1, \ldots, s_d)\), write the
value of the loss at the critical point corresponding to
\(S = \{1, \ldots, r\}\) with \(r < d\).
\end{exercise}

\begin{solution}[ex:td-1c]
At the critical point \(W = P_S M\) with
\(S = \{1, \ldots, r\}\), the student retains only the first \(r\)
modes: \(W = \text{diag}(s_1, \ldots, s_r, 0, \ldots, 0)\). The loss is:

\[\mathcal{L} = \frac{1}{2}\|M - P_S M\|_F^2 = \frac{1}{2}\sum_{\alpha \notin S} s_\alpha^2 = \frac{1}{2}\sum_{\alpha=r+1}^d s_\alpha^2\]

This is strictly positive whenever \(r < d\) (since all
\(s_\alpha > 0\)), so these are not global minima. By the Hessian
analysis (extending \cref{ex:td-1b}), the direction corresponding to
``switching on'' a missing mode \(\alpha \notin S\) is a descent
direction, making these critical points saddle points. All critical
points with \(|S| = d\) have \(W = M\) and \(\mathcal{L} = 0\): these
are the global minima. Therefore there are \textbf{no spurious local
minima} --- every local minimum is global. \(\square\)
\end{solution}

\begin{exercise}[Symmetry of the global minima]
\label{ex:td-1d}
The group
\(GL_h := GL_{d_1} \times \cdots \times GL_{d_{L-1}}\) acts on the
weights by:

\[(W_1, \ldots, W_L) \mapsto (g_1 W_1,\; g_2 W_2 g_1^{-1},\; \ldots,\; W_L g_{L-1}^{-1})\]

Verify that the student map \(\mu(\theta) = W_L \cdots W_1\) is
invariant under this action: \(\mu(g \cdot \theta) = \mu(\theta)\).
\end{exercise}

\begin{solution}[ex:td-1d]
Under the group action:

\[\mu(g \cdot \theta) = W_L g_{L-1}^{-1} \cdot g_{L-1} W_{L-1} g_{L-2}^{-1} \cdot \ldots \cdot g_1 W_1 = W_L W_{L-1} \cdots W_1 = \mu(\theta)\]

All \(g_l\) and \(g_l^{-1}\) factors cancel telescopically. This means
every parameter point in the orbit
\(\mathcal{O}_\theta = GL_h \cdot \theta\) maps to the same student
\(W\). \(\square\)

For the dimension of the set of global minima,
this invariance means the following. The global minima form the fiber
\(\mu^{-1}(M)\), which contains
the entire orbit \(GL_h \cdot \theta^*\) for any global minimizer
\(\theta^*\). The orbit has dimension \(\sum_{l=1}^{L-1} d_l^2\) (the
dimension of \(GL_h\)), so the set of global minima is a
\textbf{continuous manifold of very high dimension} --- far from being
isolated points. \(\square\)
\end{solution}

\section{Gradient flow and conserved quantities}
\label{sec:td-gradflow}

We now study gradient flow:
\(\dot{\theta}(t) = -\nabla_\theta \mathcal{L}(\theta(t))\), the
continuous-time limit of gradient descent with infinitesimal learning
rate.

\begin{exercise}[Gradient-flow equations]
\label{ex:td-2a}
For a two-layer DLN (\(L=2\)), the loss is given by \(\mathcal{L} = \frac{1}{2}\|M - W_2 W_1\|_F^2\).
Assume that each weight matrix \(W_i\) is diagonal.
Derive the gradient flow equations for each layer.
In particular, show that:
\[\dot{W}_1 = W_2^\top(M - W_2 W_1), \qquad \dot{W}_2 = (M - W_2 W_1) W_1^\top\]

\emph{Note}: these equations also hold for general (non-diagonal)
\(W_1\) and \(W_2\).
\end{exercise}

\begin{solution}[ex:td-2a]
We have
\(\mathcal{L} = \frac{1}{2}\|M - W_2 W_1\|_F^2 = \frac{1}{2}\text{Tr}\left[(M - W_2 W_1)^\top(M - W_2 W_1)\right]\).

\textbf{General matrix derivation.} Expand:

\[\mathcal{L} = \frac{1}{2}\text{Tr}(M^\top M) - \text{Tr}(M^\top W_2 W_1) + \frac{1}{2}\text{Tr}(W_1^\top W_2^\top W_2 W_1)\]

Differentiating with respect to \(W_1\), using
\(\frac{\partial}{\partial A}\text{Tr}(B^\top A) = B\) and
\(\frac{\partial}{\partial A}\text{Tr}(A^\top C A) = (C + C^\top)A\):

\[\nabla_{W_1}\mathcal{L} = -W_2^\top M + W_2^\top W_2 W_1 = -W_2^\top(M - W_2 W_1)\]

So \(\dot{W}_1 = -\nabla_{W_1}\mathcal{L} = W_2^\top(M - W_2 W_1)\).

Similarly, \(\nabla_{W_2}\mathcal{L} = -(M - W_2 W_1)W_1^\top\), giving
\(\dot{W}_2 = (M - W_2 W_1)W_1^\top\).

\textbf{Diagonal shortcut.} For diagonal matrices
\(W_1 = \text{diag}(a_\alpha)\), \(W_2 = \text{diag}(b_\alpha)\),
\(M = \text{diag}(s_\alpha)\): the loss decouples as
\(\mathcal{L} = \frac{1}{2}\sum_\alpha(s_\alpha - b_\alpha a_\alpha)^2\).
Then
\(\dot{a}_\alpha = -\partial\mathcal{L}/\partial a_\alpha = b_\alpha(s_\alpha - b_\alpha a_\alpha)\)
and \(\dot{b}_\alpha = a_\alpha(s_\alpha - b_\alpha a_\alpha)\), which
is the diagonal version of the matrix equations above. \(\square\)
\end{solution}

\begin{exercise}[Balancedness is conserved]
\label{ex:td-2b}
Define the \textbf{balancedness matrix}:

\[G := W_2^\top W_2 - W_1 W_1^\top\]

Show that \(G\) is conserved under the gradient flow,
i.e.~\(\dot{G} = 0\). In other words, gradient flow is constrained to the balanced manifold \begin{equation*}
    \mathcal{M}_{\beta}
    :=
    \{\theta\in\Omega\mid
    W_{\ell+1}^{\top}W_{\ell+1}
    -
    W_{\ell}W_{\ell}^{\top}
    =
    \beta_\ell,\;
    \ell=1,\ldots,L-1\}.
\end{equation*}

\emph{Hint}: Compute
\(\dot{G}\),
substitute the gradient flow equations and verify that the
terms cancel pairwise.
\end{exercise}

\begin{solution}[ex:td-2b]
Let \(E := M - W_2 W_1\) denote the residual. Compute:

\[\dot{G} = \dot{W}_2^\top W_2 + W_2^\top \dot{W}_2 - \dot{W}_1 W_1^\top - W_1 \dot{W}_1^\top\]

Substitute the gradient flow equations. Using \(\dot{W}_2 = E W_1^\top\)
and \(\dot{W}_1 = W_2^\top E\):

\[\dot{W}_2^\top W_2 = (E W_1^\top)^\top W_2 = W_1 E^\top W_2\]

\[W_2^\top \dot{W}_2 = W_2^\top E W_1^\top\]

\[\dot{W}_1 W_1^\top = W_2^\top E W_1^\top\]

\[W_1 \dot{W}_1^\top = W_1(W_2^\top E)^\top = W_1 E^\top W_2\]

Therefore:

\[\dot{G} = W_1 E^\top W_2 + W_2^\top E W_1^\top - W_2^\top E W_1^\top - W_1 E^\top W_2 = 0\]

The terms cancel pairwise. \(G\) is conserved. \(\square\)
\end{solution}

From now on, assume \textbf{balanced initialization}:
\(G(0) = 0\), which by \cref{ex:td-2b} means \(W_2^\top W_2 = W_1 W_1^\top\)
for all time. We want to derive the gradient flow in function space (the
ODE for the student \(W = W_2 W_1\)). This requires several steps.

\begin{exercise}[Function-space velocity]
\label{ex:td-2ci}
Compute \(\dot{W} := \frac{d}{dt}(W_2 W_1)\) using the
gradient flow equations from \cref{ex:td-2a}. Show that:

\[\dot{W} = (M - W)\, W_1^\top W_1 + W_2 W_2^\top\, (M - W)\]
\end{exercise}

\begin{solution}[ex:td-2ci]
Apply the product rule:

\[\dot{W} = \dot{W}_2 W_1 + W_2 \dot{W}_1 = (M - W_2 W_1)W_1^\top \cdot W_1 + W_2 \cdot W_2^\top(M - W_2 W_1)\]

Writing \(W = W_2 W_1\):

\[\dot{W} = (M - W)\,W_1^\top W_1 + W_2 W_2^\top\,(M - W) \qquad \square\]
\end{solution}

\begin{exercise}[\(W_1^\top W_1\) in terms of \(W\)]
\label{ex:td-2cii}
We now need to express \(W_1^\top W_1\) and
\(W_2 W_2^\top\) in terms of \(W = W_2 W_1\). Show that:

\[W_1^\top W_1 = (W^\top W)^{1/2}\]
\end{exercise}

\begin{solution}[ex:td-2cii]
Balancedness gives \(W_2^\top W_2 = W_1 W_1^\top\), so

\[W^\top W = W_1^\top \left(W_2^\top W_2\right) W_1 = W_1^\top \left(W_1 W_1^\top\right) W_1 = \left(W_1^\top W_1\right)^2\]

Both \(W^\top W\) and \(W_1^\top W_1\) are positive semidefinite.
A positive semidefinite matrix has a unique positive semidefinite square root, so it follows that
\[(W^\top W)^{1/2} = W_1^\top W_1 \qquad \square\]
\end{solution}

\begin{exercise}[\(W_2 W_2^\top\) in terms of \(W\)]
\label{ex:td-2ciii}
Similarly, show that
\(W_2 W_2^\top = (W W^\top)^{1/2}\).
\end{exercise}

\begin{solution}[ex:td-2ciii]
Analogous to \cref{ex:td-2cii}.
\end{solution}

\begin{exercise}[The balanced function-space ODE]
\label{ex:td-2civ}
Substitute the results of \cref{ex:td-2cii,ex:td-2ciii} into \cref{ex:td-2ci}
to obtain:

\[\dot{W} = (W W^\top)^{1/2}(M - W) + (M - W)(W^\top W)^{1/2}\]
\end{exercise}

\begin{solution}[ex:td-2civ]
Substituting into \cref{ex:td-2ci}:

\[\dot{W} = (M - W)(W^\top W)^{1/2} + (WW^\top)^{1/2}(M - W) \qquad \square\]
\end{solution}

\begin{exercise}[The NTK operator]
\label{ex:td-2d}
Define the NTK operator for \(L = 2\) as:

\[K[F] := (WW^\top)^{1/2}\, F + F\, (W^\top W)^{1/2}\]

Show that the gradient flow from \cref{ex:td-2civ} can be written as
\(\dot{W} = K[M - W]\).

It turns out that this result generalizes to depth \(L\) on the balanced
manifold:

\[\dot{W} = \sum_{k=1}^{L} (WW^\top)^{\frac{L-k}{L}} (M - W) (W^\top W)^{\frac{k-1}{L}}\]
This is the NTK equation in the case of DLNs. The NTK equation is a gradient flow in function space with NTK being a preconditioning operator for the gradient.
\end{exercise}

\begin{solution}[ex:td-2d]
With \(F = M - W\), the equation from \cref{ex:td-2civ} reads:

\[\dot{W} = (WW^\top)^{1/2}\,F + F\,(W^\top W)^{1/2} = K[F] = K[M - W] \qquad \square\]

The general depth-\(L\) result
\(\dot{W} = \sum_{k=1}^{L}(WW^\top)^{(L-k)/L}(M-W)(W^\top W)^{(k-1)/L}\)
follows from the same approach applied to the \(L\)-fold balanced
conditions \(W_{l+1}^\top W_{l+1} = W_l W_l^\top\) for all \(l\).
\end{solution}

\section{Rich regime: saddle-to-saddle}
\label{sec:td-rich}

This is the core problem. We derive the exact solution of the rich
regime dynamics directly from the self-consistent equation of \cref{sec:td-gradflow},
emphasizing the NTK perspective.

\textbf{Alignment assumption.} Start from the balanced
gradient flow equation derived in \cref{ex:td-2civ}:

\[\dot{W} = (WW^\top)^{1/2}(M - W) + (M - W)(W^\top W)^{1/2}\]

We work in the \textbf{rich regime} (small initialization) with balanced
weights. Assume that the student is \textbf{aligned} to the teacher: the
singular vectors of \(W(t)\) coincide with those of \(M\) at all times.
Concretely, write:

\[M = U\,\text{diag}(s_1, \ldots, s_d)\,V^\top, \qquad W(t) = U\,\text{diag}(w_1(t), \ldots, w_d(t))\,V^\top\]

where \(U, V\) are the left/right singular vectors of \(M\), and
\(w_\alpha(t) \geq 0\) are the evolving singular values of \(W\).

\begin{exercise}[Aligned NTK]
\label{ex:td-3ai}
Under this alignment assumption, show that the NTK
operator is aligned to the task, i.e.~\(K[F]\) preserves the SVD basis.
In particular, show that:

\begin{align*}
(WW^\top)^{1/2} &= U\,\text{diag}(w_1, \ldots, w_d)\,U^\top, \\
(W^\top W)^{1/2} &= V\,\text{diag}(w_1, \ldots, w_d)\,V^\top
\end{align*}

and that the residual is
\(M - W = U\,\text{diag}(s_1 - w_1, \ldots, s_d - w_d)\,V^\top\).

\emph{Hint}: Recall that for \(A = U\,\text{diag}(\sigma_i)\,V^\top\),
we have \(AA^\top = U\,\text{diag}(\sigma_i^2)\,U^\top\), and therefore
\((AA^\top)^{1/2} = U\,\text{diag}(|\sigma_i|)\,U^\top\).
\end{exercise}

\begin{solution}[ex:td-3ai]
Under the alignment assumption,
\(W = U\,\text{diag}(w_\alpha)\,V^\top\). Then:

\[WW^\top = U\,\text{diag}(w_\alpha^2)\,U^\top\]

Since \(w_\alpha \geq 0\), the positive matrix square root is:

\[(WW^\top)^{1/2} = U\,\text{diag}(w_\alpha)\,U^\top\]

Similarly:

\[W^\top W = V\,\text{diag}(w_\alpha^2)\,V^\top \quad \Longrightarrow \quad (W^\top W)^{1/2} = V\,\text{diag}(w_\alpha)\,V^\top\]

The residual is immediate:

\[M - W = U\,\text{diag}(s_\alpha - w_\alpha)\,V^\top\]

The NTK operator acts on any matrix
\(F = U\,\text{diag}(f_\alpha)\,V^\top\) as:

\[
\begin{aligned}
K[F]
&=
U\,\text{diag}(w_\alpha)\,U^\top
U\,\text{diag}(f_\alpha)\,V^\top \\
&\quad
+ U\,\text{diag}(f_\alpha)\,V^\top
V\,\text{diag}(w_\alpha)\,V^\top \\
&=
U\,\text{diag}(2w_\alpha f_\alpha)\,V^\top.
\end{aligned}
\]

So \(K[F]\) stays in the SVD basis of \(M\) --- the NTK is aligned to
the task. \(\square\)
\end{solution}

\begin{exercise}[Decoupled scalar ODEs]
\label{ex:td-3aii}
Substitute into the self-consistent equation. Show that
the matrix equation decouples into \(d\) independent scalar ODEs:

\[\dot{w}_\alpha = 2\,w_\alpha\,(s_\alpha - w_\alpha), \qquad \alpha = 1, \ldots, d\]

\emph{Hint}: Compute the product \((WW^\top)^{1/2}(M - W)\) in the SVD
basis. It is diagonal with entries \(w_\alpha(s_\alpha - w_\alpha)\).
The second term contributes identically, giving the factor of 2.
\end{exercise}

\begin{solution}[ex:td-3aii]
Substitute into
\(\dot{W} = (WW^\top)^{1/2}(M-W) + (M-W)(W^\top W)^{1/2}\):

\textbf{First term:}

\[U\,\text{diag}(w_\alpha)\,\underbrace{U^\top U}_{I}\,\text{diag}(s_\alpha - w_\alpha)\,V^\top = U\,\text{diag}\big(w_\alpha(s_\alpha - w_\alpha)\big)\,V^\top\]

\textbf{Second term:}

\[U\,\text{diag}(s_\alpha - w_\alpha)\,\underbrace{V^\top V}_{I}\,\text{diag}(w_\alpha)\,V^\top = U\,\text{diag}\big((s_\alpha - w_\alpha)\,w_\alpha\big)\,V^\top\]

Summing:

\[\dot{W} = U\,\text{diag}\big(2\,w_\alpha(s_\alpha - w_\alpha)\big)\,V^\top\]

Since \(W = U\,\text{diag}(w_\alpha)\,V^\top\) and the singular vectors
are constant by assumption, we read off:

\[\boxed{\dot{w}_\alpha = 2\,w_\alpha\,(s_\alpha - w_\alpha)}\]

The NTK perspective makes the structure transparent: the factor
\(2w_\alpha\) is the NTK eigenvalue for mode \(\alpha\), which amplifies
learning in directions that are already strong, while
\((s_\alpha - w_\alpha)\) is the residual. \(\square\)
\end{solution}

\begin{exercise}[Timescale of learning]
\label{ex:td-3b}
The ODE
\(\dot{w} = 2w(s - w)\) is a logistic equation whose solution is:

\[w(t) = \frac{s}{1 + \left(\frac{s}{w_0} - 1\right)e^{-2st}}\]

where \(w_0 = w(0)\). Compute the time \(t_\alpha\) it takes for mode
\(\alpha\) to travel from initial strength \(w_0\) to a final strength
\(w_f\). Show that:

\[t_\alpha = \frac{1}{2 s_\alpha} \ln\left(\frac{w_f\,(s_\alpha - w_0)}{w_0\,(s_\alpha - w_f)}\right)\]

Deduce that modes with larger singular values \(s_\alpha\) are learned
faster. This is a \textbf{separation of timescales}: the network learns
features in decreasing order of their singular value strength.
\end{exercise}

\begin{solution}[ex:td-3b]
From the logistic solution
\(w(t) = \frac{s}{1 + (s/w_0 - 1)e^{-2st}}\), we invert to find \(t\) as
a function of \(w\). The derivation is equivalent to integrating
\(\dot{w} = 2w(s-w)\) by separation of variables:

\[\int_{w_0}^{w_f} \frac{dw}{w(s-w)} = 2\,t_\alpha\]

Using partial fractions
\(\frac{1}{w(s-w)} = \frac{1}{s}\!\left(\frac{1}{w} + \frac{1}{s-w}\right)\):

\[\frac{1}{s_\alpha}\Big[\ln w - \ln(s_\alpha - w)\Big]_{w_0}^{w_f} = 2\,t_\alpha\]

\[\frac{1}{s_\alpha}\ln\frac{w_f\,(s_\alpha - w_0)}{w_0\,(s_\alpha - w_f)} = 2\,t_\alpha\]

Therefore:

\[\boxed{t_\alpha = \frac{1}{2s_\alpha}\ln\left(\frac{w_f\,(s_\alpha - w_0)}{w_0\,(s_\alpha - w_f)}\right)}\]

For a fixed ratio \(w_f/s_\alpha\) and fixed \(w_0\), the learning time
scales as \(t_\alpha \propto 1/s_\alpha\): \textbf{modes with larger
singular values are learned faster}. The network learns features in
decreasing order of their strength --- a strong separation of
timescales.

\textbf{Contrast with the lazy regime} (anticipating \cref{sec:td-lazy}): in the
lazy regime the NTK is frozen at \(K_0 = 2w_0 \gg 1\), so
\(\dot{w}_\alpha \approx 2w_0(s_\alpha - w_\alpha)\). All modes converge
at the same exponential rate \(2w_0\), independent of \(s_\alpha\). The
rich regime has state-dependent NTK (\(2w_\alpha\)), creating a positive
feedback loop --- modes that are already large learn even faster ---
which amplifies the differences between \(s_\alpha\) into a hierarchy of
timescales \(t_1 \ll t_2 \ll \cdots \ll t_r\). \(\square\)
\end{solution}

\begin{exercise}[Incremental learning]
\label{ex:td-3c}
Consider a teacher with
\(r\) nonzero singular values and a small uniform initialization
\(w_\alpha(0) = w_0 \ll s_r\) for all \(\alpha\). Qualitatively, discuss:

\begin{itemize}
\tightlist
\item
  What is the approximate distribution of singular values of \(W(t)\) at early times?
\item
  How do the singular values of $W(t)$ evolve throughout training?
\item
  What happens in the limit \(t \to \infty\)?
\end{itemize}

What does this analysis tell us about the sequence of critical points visited by
the gradient flow?
\end{exercise}

\begin{solution}[ex:td-3c]
With uniform initialization \(w_\alpha(0) = w_0\) for all
\(\alpha\), the sigmoid solution shows that mode \(\alpha\) remains near
\(w_0\) until time \(t \approx \frac{1}{2s_\alpha}\ln(s_\alpha/w_0)\),
then transitions rapidly to \(w_\alpha \approx s_\alpha\).

\textbf{Early times} (\(t \sim \frac{1}{2s_1}\ln(s_1/w_0)\)): Only mode
1 has grown significantly, so the spectrum is \textbf{one singular value
near \(s_1\), with all the others still near \(w_0 \ll s_r\)}. The
student is approximately \(W(t) \approx w_1(t)\,u_1 v_1^\top\),
effectively \textbf{rank 1}.

\textbf{Intermediate times}: As \(t\) increases, \(w_2\) switches on,
then \(w_3\), etc. The singular values \textbf{switch on one at a time},
in decreasing order of \(s_\alpha\), each staying near \(w_0\) and then
rising to \(s_\alpha\) over a narrow window. The student is
approximately
\(W(t) \approx \sum_{\alpha=1}^{k(t)} w_\alpha(t)\,u_\alpha v_\alpha^\top\)
where \(k(t)\) increases stepwise, so the effective \textbf{rank
increases incrementally}.

\textbf{Long times} (\(t \to \infty\)): All modes converge,
\(w_\alpha \to s_\alpha\), and \(W \to M\).

This is an \textbf{implicit bias toward low-rank (simple) solutions}: at
any finite time, the network represents the best rank-\(k\)
approximation to the teacher. The gradient flow effectively performs a
greedy SVD.

\textbf{Connection to the saddle structure (\cref{ex:td-1c}):} The saddle
points \(P_S M\) with \(|S| = k\) are exactly the best rank-\(k\)
approximations to \(M\), with loss
\(\frac{1}{2}\sum_{\alpha > k} s_\alpha^2\). The gradient flow
trajectory passes \textbf{near} these saddles as it incrementally
recruits modes. The \textbf{plateaus} in the loss curve correspond to
time spent near saddle points (only \(k\) modes active), and the
\textbf{sharp transitions} correspond to escaping along the unstable
direction that activates the next mode. \(\square\)
\end{solution}

\section{Lazy regime}
\label{sec:td-lazy}

We now show that large initialization freezes the NTK and eliminates the
timescale separation found in \cref{sec:td-rich}. For clarity, we work with the
\textbf{diagonal scalar} model from \cref{ex:td-1a}, so each mode
\(\alpha\) is independent. This avoids matrix algebra and isolates the
essential mechanism.

\textbf{Setup.} Consider a single mode: a depth-2 diagonal
DLN with scalar weights \(a, b\) learning a target \(s > 0\). From
\cref{sec:td-gradflow,sec:td-rich}, the balanced gradient flow for \(w = ab\) is:

\[\dot{w} = 2w(s - w)\]

The factor \(2w\) is the (scalar) NTK --- it is
\textbf{state-dependent}: the effective learning rate depends on the
current value of \(w\).

\begin{exercise}[Linearizing the NTK]
\label{ex:td-4a}
Now initialize at a \textbf{large} value \(w_0 \gg s > 0\). Assume that in
the early phase of training, while \(w\) has not changed much from
\(w_0\), the ODE becomes approximately:

\[\dot{w} \approx 2w_0(s - w)\]
\end{exercise}

\begin{solution}[ex:td-4a]
Starting from the same ODE \(\dot{w} = 2w(s - w)\), with
\(w_0 \gg s > 0\), the weight \(w\) needs to decrease from \(w_0\) to
\(s\). As long as \(w\) has not changed much from \(w_0\), we can write
\(w = w_0 + \delta w\) with \(|\delta w| \ll w_0\), so:

\[2w = 2(w_0 + \delta w) \approx 2w_0\]

The ODE becomes:

\[\dot{w} \approx 2w_0(s - w) \qquad \square\]

This is a \textbf{linear} ODE --- the NTK factor \(2w\) has been frozen
at its initial value \(2w_0\). Note that only the NTK prefactor is
linearized; the residual \((s - w)\) is kept exact since it drives
learning.
\end{solution}

\begin{exercise}[Frozen-NTK solution]
\label{ex:td-4b}
Solve the linearized ODE from \cref{ex:td-4a}. Show that:

\[w(t) = s + (w_0 - s)\,e^{-2w_0 t}\]

What is the learning timescale?
\end{exercise}

\begin{solution}[ex:td-4b]
The linearized ODE \(\dot{w} = 2w_0(s - w)\) is first-order
linear. Set \(\tilde{w} := w - s\), then
\(\dot{\tilde{w}} = -2w_0\,\tilde{w}\), with solution
\(\tilde{w}(t) = (w_0 - s)\,e^{-2w_0 t}\). Therefore:

\[\boxed{w(t) = s + (w_0 - s)\,e^{-2w_0 t}}\]

This is \textbf{exponential} convergence to \(s\) (contrast with the
sigmoid of \cref{sec:td-rich}). The learning timescale is:

\[t_{\text{lazy}} \sim \frac{1}{2w_0}\]

It depends on the \textbf{initialization scale} \(w_0\) but \textbf{not
on the target} \(s\). This is the defining feature of the lazy regime:
the learning speed is set by the NTK at initialization, not by the
structure of the task.

\textbf{Self-consistency:} During the learning phase
(\(t \lesssim 1/(2w_0)\)), the displacement is
\(|w(t) - w_0| \leq |s - w_0| \approx w_0\) (since \(w_0 \gg s\)). The
relative change in the NTK is
\(\Delta(2w)/(2w_0) = (w_0 - s)/w_0 = 1 - s/w_0 \approx 1\), so the NTK
does change significantly in absolute terms. However, the key point is
that the dynamics remain well approximated by the linear ODE because the
convergence rate is dominated by \(2w_0\), which is large and
approximately constant throughout training. In the rich regime, by
contrast, the NTK changes from \(2w_0 \approx 0\) to \(2s\) --- a change
of order \(s/w_0 \to \infty\) relative to the initial value --- making
the linearization completely invalid. \(\square\)
\end{solution}

\begin{exercise}[No timescale separation]
\label{ex:td-4c}
Now restore the mode
index \(\alpha\). In the lazy regime, each mode satisfies
\(\dot{w}_\alpha \approx 2w_0(s_\alpha - w_\alpha)\) with the
\textbf{same} rate \(2w_0\) for all \(\alpha\). Compare the lazy
learning time \(t_\alpha^{\text{lazy}}\) with the rich-regime timescale
from \cref{ex:td-3b}: \(t_\alpha^{\text{rich}} \propto 1/s_\alpha\).
\end{exercise}

\begin{solution}[ex:td-4c]
Restoring the mode index, each mode satisfies
\(\dot{w}_\alpha \approx 2w_0(s_\alpha - w_\alpha)\) with solution:

\[w_\alpha(t) = s_\alpha + (w_0 - s_\alpha)\,e^{-2w_0 t}\]

All modes converge at the \textbf{same exponential rate} \(2w_0\),
regardless of \(s_\alpha\). The time for mode \(\alpha\) to go from
\(w_0\) to within \((1-\epsilon)\) of \(s_\alpha\) is:

\[t_\alpha^{\text{lazy}} = \frac{1}{2w_0}\ln\!\left(\frac{w_0 - s_\alpha}{\epsilon\,(w_0 - s_\alpha)}\right) = \frac{1}{2w_0}\ln\frac{1}{\epsilon}\]

which is \textbf{independent of \(s_\alpha\)}. All modes reach their
targets simultaneously.

\textbf{Contrast with the rich regime:}

\begin{center}
\begin{tabular}{p{0.30\linewidth}p{0.27\linewidth}p{0.27\linewidth}}
\hline
 & Rich regime & Lazy regime \\
\hline
NTK & State-dependent: \(2w_\alpha\) & Frozen: \(2w_0\) \\
Learning time for mode \(\alpha\) & \(t_\alpha \propto 1/s_\alpha\) & \(t_\alpha \approx \text{const}\) \\
Timescale separation & Strong: \(t_1 \ll t_2 \ll \cdots\) & None \\
Intermediate solutions & Low-rank (incremental) & Full-rank from the start \\
Implicit rank bias & Yes (low-rank → high-rank) & No \\
\hline
\end{tabular}
\end{center}
In the lazy regime, all components of \(M\) are learned simultaneously:
signal and noise alike. The network converges directly to the full OLS
solution without passing through low-rank intermediates. There is no
mechanism to separate signal from noise based on singular value
structure. With early stopping, the rich regime recovers a low-rank
signal while ignoring noise; the lazy regime cannot. \(\square\)
\end{solution}

\section{Mixed dynamics: unifying lazy and rich}
\label{sec:td-mixed}

In practice, networks are neither purely lazy nor purely rich. Tu,
Aranguri \& Jacot (2024) provide a unified description. We develop the
key ideas using the diagonal scalar model.

\begin{exercise}[The interpolating ODE]
\label{ex:td-5a}
In \cref{sec:td-rich,sec:td-lazy}, we
studied the same ODE \(\dot{w}_\alpha = 2w_\alpha(s_\alpha - w_\alpha)\)
in two limits: small \(w_0\) (rich) and large \(w_0\) (lazy). A more
general model for a depth-2 DLN with balanced initialization at scale
\(\sigma\) and width \(n\) replaces the scalar NTK \(2w_\alpha\) by:

\[\dot{w}_\alpha = 2\sqrt{w_\alpha^2 + \tau^2}\;(s_\alpha - w_\alpha)\]

where \(\tau := \sigma^2 \sqrt{n}\) is a \textbf{threshold} parameter
that depends on the initialization scale and width.

Verify that this ODE reproduces the two known regimes:

\begin{itemize}
\tightlist
\item
  \textbf{Rich limit} (\(\tau \to 0\)): recover
  \(\dot{w}_\alpha = 2|w_\alpha|(s_\alpha - w_\alpha)\).
\item
  \textbf{Lazy limit} (\(\tau \to \infty\) with \(w_\alpha\) bounded):
  recover \(\dot{w}_\alpha \approx 2\tau(s_\alpha - w_\alpha)\).
\end{itemize}
\end{exercise}

\begin{solution}[ex:td-5a]
The interpolating ODE is
\(\dot{w}_\alpha = 2\sqrt{w_\alpha^2 + \tau^2}\,(s_\alpha - w_\alpha)\).

\textbf{Rich limit (\(\tau \to 0\)):} The square root reduces to
\(\sqrt{w_\alpha^2} = |w_\alpha|\), giving:

\[\dot{w}_\alpha = 2|w_\alpha|(s_\alpha - w_\alpha)\]

For \(w_\alpha > 0\) this is \(2w_\alpha(s_\alpha - w_\alpha)\), exactly
the logistic ODE from \cref{ex:td-3aii}. \(\square\)

\textbf{Lazy limit (\(\tau \to \infty\), \(w_\alpha\) bounded):} The
square root is dominated by \(\tau\):
\(\sqrt{w_\alpha^2 + \tau^2} \approx \tau\), giving:

\[\dot{w}_\alpha \approx 2\tau(s_\alpha - w_\alpha)\]

This is a linear ODE with rate \(2\tau\), independent of \(s_\alpha\)
--- the frozen-NTK regime of \cref{sec:td-lazy} (with \(\tau\) playing the role
of \(w_0\)). \(\square\)
\end{solution}

\begin{exercise}[Two-phase dynamics]
\label{ex:td-5b}
At initialization, all modes
start at \(w_\alpha(0) \sim \sigma^2 \ll \tau\) (since
\(\sqrt{n} \gg 1\)). Observe that:

\begin{itemize}
\tightlist
\item
  When \(|w_\alpha| \ll \tau\), the effective learning rate is
  \(\approx 2\tau\), the same for all modes (lazy behavior).
\item
  When \(|w_\alpha| \gg \tau\), the effective learning rate is
  \(\approx 2|w_\alpha|\), which is mode-dependent (rich behavior).
\end{itemize}

Describe the resulting two-phase dynamics in words: what happens first,
and what happens later?
\end{exercise}

\begin{solution}[ex:td-5b]
At initialization,
\(w_\alpha(0) \sim \sigma^2 \ll \tau = \sigma^2\sqrt{n}\) (since
\(\sqrt{n} \gg 1\)). So initially all modes satisfy
\(|w_\alpha| \ll \tau\).

\textbf{Early phase (lazy):} \(\sqrt{w_\alpha^2 + \tau^2} \approx \tau\)
for all \(\alpha\). Every mode evolves at the same linear rate
\(2\tau\):

\[\dot{w}_\alpha \approx 2\tau(s_\alpha - w_\alpha) \qquad \Longrightarrow \qquad w_\alpha(t) \approx s_\alpha(1 - e^{-2\tau t})\]

(assuming \(w_\alpha(0) \approx 0\)). The dynamics are approximately
linear and the NTK is approximately constant. During this phase, the
network \textbf{aligns} with the task: each \(w_\alpha\) grows toward
\(s_\alpha\) at the same rate. There is no timescale separation.

\textbf{Late phase (rich):} Once some \(w_\alpha\) grow past \(\tau\),
the square root transitions to
\(\sqrt{w_\alpha^2 + \tau^2} \approx |w_\alpha|\), and those modes enter
the rich regime with the sigmoidal, self-accelerating dynamics
\(\dot{w}_\alpha \approx 2w_\alpha(s_\alpha - w_\alpha)\). The
state-dependent NTK creates timescale separation, and the modes that
have crossed the threshold converge rapidly to their targets.

In summary: the network starts in a lazy phase where all modes grow
uniformly (alignment), then transitions to a rich phase where modes
accelerate individually (incremental learning). \(\square\)
\end{solution}

\begin{exercise}[Mode-by-mode transition]
\label{ex:td-5c}
Not all modes cross the
threshold \(\tau\) at the same time: which modes cross it first? Explain
why this means that a single network can simultaneously have some modes
in the rich regime and others still in the lazy regime.
\end{exercise}

\begin{solution}[ex:td-5c]
In the lazy phase, each mode grows as
\(w_\alpha(t) \approx s_\alpha(1 - e^{-2\tau t})\). At any given time,
\(w_\alpha(t) \propto s_\alpha\): \textbf{modes with larger \(s_\alpha\)
are larger}. Since the lazy-to-rich transition occurs when
\(|w_\alpha| \sim \tau\), the time for mode \(\alpha\) to cross the
threshold satisfies:

\[s_\alpha(1 - e^{-2\tau t_\alpha^*}) = \tau \qquad \Longrightarrow \qquad t_\alpha^* = \frac{1}{2\tau}\ln\frac{s_\alpha}{s_\alpha - \tau}\]

For \(s_\alpha > \tau\), this crossing time exists and is smaller for
larger \(s_\alpha\). For \(s_\alpha < \tau\), the mode never reaches the
threshold and remains permanently lazy.

This means that \textbf{different modes can be in different regimes
simultaneously}: large-\(s_\alpha\) modes have crossed into the rich
regime and are converging rapidly, while small-\(s_\alpha\) modes are
still in the lazy phase (or may never leave it). The network
interpolates between the two regimes mode by mode.

\textbf{Connection to grokking:} This lazy-to-rich transition provides a
mechanism for grokking. During the lazy phase, the network fits the
training data via kernel regression with the initial (misaligned) NTK
--- it \textbf{memorizes}. The test loss plateaus because the frozen
kernel cannot capture the task structure. Later, as modes cross the
threshold into the rich regime, feature learning kicks in: the NTK
rotates to align with the task, and the test loss drops ---
\textbf{generalization} is achieved. The grokking delay is controlled by
the time it takes the relevant modes to cross \(\tau\). Increasing the
initialization scale \(\sigma\) or the width \(n\) increases \(\tau\),
which extends the lazy phase and widens the grokking gap. This is the
same mechanism identified by Kumar et al., \emph{Grokking as the
Transition from Lazy to Rich Training Dynamics} (arXiv:2310.06110).
\(\square\)
\end{solution}

\section{Stochastic implicit bias (bonus)}
\label{sec:td-stochastic}

SGD introduces noise from mini-batching. A continuous model is the Langevin SDE:

\[d\theta_t = -\nabla \mathcal{L}(\theta_t)\,dt + \sqrt{\eta\,\Sigma(\theta_t)}\,dB_t\]

where \(\Sigma(\theta)\) is the covariance of the stochastic gradient
noise and \(\eta\) is the learning rate.

\begin{exercise}[Boltzmann equilibrium]
\label{ex:td-6a}
The probability density \(p(\theta, t)\) of the parameters
evolves according to the Fokker--Planck equation:

\[\partial_t p = -\nabla \cdot \mathbf{j}, \qquad \mathbf{j} = -\nabla \mathcal{L}(\theta)\,p(\theta) - \frac{\eta}{2}\nabla \cdot \left(\Sigma(\theta)\,p(\theta)\right)\]

where \(\mathbf{j}\) is the probability current. Assume: (i)
stationarity \(\partial_t p^* = 0\), (ii) thermal equilibrium
\(\mathbf{j} = 0\), and (iii) isotropic noise \(\Sigma = \sigma^2 I\).
Show that the equilibrium distribution is the Boltzmann distribution:

\[p^*(\theta) \propto \exp\left(-\frac{2}{\eta \sigma^2}\mathcal{L}(\theta)\right)\]

\emph{Hint}: Setting \(\mathbf{j} = 0\) with \(\Sigma = \sigma^2 I\)
gives \(\nabla \mathcal{L}\, p + \frac{\eta\sigma^2}{2}\nabla p = 0\).
This is a first-order ODE for \(p\) in terms of \(\mathcal{L}\). Try the
ansatz \(p \propto e^{-\beta \mathcal{L}}\) and solve for \(\beta\).
\end{exercise}

\begin{solution}[ex:td-6a]
Setting \(\mathbf{j} = 0\) (thermal equilibrium) with
isotropic noise \(\Sigma = \sigma^2 I\):

\[0 = -\nabla\mathcal{L}(\theta)\,p^*(\theta) - \frac{\eta\sigma^2}{2}\nabla p^*(\theta)\]

Rearranging:

\[\frac{\nabla p^*}{p^*} = -\frac{2}{\eta\sigma^2}\nabla\mathcal{L}\]

The left-hand side is \(\nabla\ln p^*\), so:

\[\nabla\ln p^*(\theta) = -\frac{2}{\eta\sigma^2}\nabla\mathcal{L}(\theta)\]

Integrating both sides:

\[\ln p^*(\theta) = -\frac{2}{\eta\sigma^2}\mathcal{L}(\theta) + \text{const}\]

\[\boxed{p^*(\theta) \propto \exp\!\left(-\frac{2}{\eta\sigma^2}\mathcal{L}(\theta)\right)}\]

This is the Boltzmann distribution with inverse temperature
\(\beta = 2/(\eta\sigma^2)\). \(\square\)
\end{solution}

\begin{exercise}[Temperature and flatness]
\label{ex:td-6b}
The ratio \(\frac{2}{\eta \sigma^2}\) plays the role of an
inverse temperature \(\beta\). Interpret what happens to the equilibrium
distribution when:

\begin{itemize}
\tightlist
\item
  \(\eta\) is very small (low temperature)
\item
  \(\eta\) is very large (high temperature)
\end{itemize}

Which regime favors flatter minima, and why might this be beneficial for
generalization?
\end{exercise}

\begin{solution}[ex:td-6b]
The effective temperature is \(T = \eta\sigma^2/2\).

\textbf{Small \(\eta\) (low temperature, \(\beta \to \infty\)):} The
distribution concentrates sharply on the global minima of
\(\mathcal{L}\). SGD converges to the lowest-loss minimizer without
exploring.

\textbf{Large \(\eta\) (high temperature, \(\beta \to 0\)):} The
distribution becomes nearly uniform over parameter space. SGD explores
broadly and does not settle into any particular minimum.

\textbf{Flat vs.~sharp minima:} At moderate temperature, the Boltzmann
distribution assigns more probability mass to \textbf{broad basins}
(flat minima) than to narrow ones. This is because a flat minimum
occupies a larger volume of parameter space at any given loss level ---
the width of the basin acts as an entropic contribution. Flat minima
tend to generalize better because small perturbations to the parameters
(or slight distribution shift in the data) do not dramatically change
the loss. Thus the implicit bias of SGD noise toward flat minima is
beneficial for generalization. \(\square\)
\end{solution}

\begin{exercise}[Anisotropic noise]
\label{ex:td-6c}
In practice, SGD noise is \textbf{not} isotropic:
\(\Sigma(\theta)\) depends on both the loss landscape and the data.
Without solving anything, explain qualitatively why anisotropic noise
can introduce an implicit bias that goes beyond what the loss function
\(\mathcal{L}\) alone would select. Specifically, why might SGD
preferentially escape sharp directions of the loss while remaining
stable along flat directions?
\end{exercise}

\begin{solution}[ex:td-6c]
When \(\Sigma(\theta)\) is anisotropic, the noise strength
varies by direction. The Fokker-Planck current becomes:

\[\mathbf{j} = -\nabla\mathcal{L}\,p - \frac{\eta}{2}\nabla\cdot(\Sigma(\theta)\,p)\]

The second term introduces an \textbf{Itô drift}
\(\frac{\eta}{2}\nabla\cdot\Sigma(\theta)\) that depends on the spatial
variation of the noise covariance. In directions where \(\Sigma\) has
large eigenvalues (high noise), the effective diffusion is strong and
the system escapes easily from sharp regions of the loss landscape. In
directions where \(\Sigma\) has small eigenvalues (low noise), the
system is more stable and tends to remain there.

This creates a \textbf{direction-dependent regularizer}: SGD
preferentially pushes parameters out of sharp directions (high gradient
variance → large noise eigenvalue → fast escape) while preserving
parameters along flat directions (low gradient variance → small noise
eigenvalue → stability). This goes beyond what the Boltzmann
distribution on \(\mathcal{L}\) alone would predict. In general,
detailed balance \(\mathbf{j} = 0\) may not hold for anisotropic,
state-dependent noise, leading to persistent probability currents and
non-equilibrium steady states that further modify the implicit bias.
\(\square\)
\end{solution}

\section*{Learn more}

The readings are roughly ordered in a way that makes sense for learning. Discuss the readings in groups of 2 or 3, to be formed in this \href{https://docs.google.com/document/d/13u-zhLUnAT9Wzn2IUsbNUOt9vNE5u8QV8GV1NAr1gow/edit?tab=t.0}{Google doc}. Spend 30 minutes reading, 30 minutes discussing and 30 minutes to write down your thoughts in the Google doc.

\subsection*{Key readings}
\begin{itemize}
  \item The strong default suggestion for everyone is: Saxe et al. \href{https://www.pnas.org/doi/10.1073/pnas.1820226116}{A mathematical theory of semantic development in deep neural networks} --- Read the supplementary materials section to understand the exact solution of gradient flow dynamics in deep linear networks
  \item For people with a strong algebraic geometry and representation geometry background read: \href{https://arxiv.org/abs/2411.19920}{Geometry of fibers of the multiplication map of deep linear neural networks} Simon Pepin Lehalleur et al. --- Stratifies global minima into orbits using quiver representation theory
  \item For people with a strong dynamical system and differential geometry background read: \href{https://arxiv.org/abs/2411.09004}{The geometry of the deep linear network} Govind Menon --- Beautiful mathematical treatment of gradient flow in DLNs and a surprising mathematical result relating gradient flow and free-energy minimization. Makes implicit regularization explicit (under balanced assumptions)
  \item For people who are much more empirically minded: \href{https://arxiv.org/abs/2602.07852}{Emergent Misalignment is Easy, Narrow Misalignment is Hard} --- When fine-tuned on data with harmful inputs, the model generalizes in harmful ways on other unrelated datasets. By regularizing, we can mitigate emergent misalignment
  \item Literature Review: \href{https://arxiv.org/abs/2604.21691}{There Will Be a Scientific Theory of Deep Learning}
  \item To go further: \href{https://arxiv.org/abs/2506.06489}{Alternating Gradient Flows: A Theory of Feature Learning in Two-layer Neural Networks}
\end{itemize}

\subsection*{Loss landscape geometry}
\begin{itemize}
  \item \href{https://arxiv.org/abs/1605.07110}{Deep Learning without Poor Local Minima}, Kenji Kawaguchi --- For non-bottlenecked DLNs, all local minima are global
  \item \href{https://arxiv.org/abs/2107.13289}{The loss landscape of deep linear neural networks: a second-order analysis} Achour et al. --- First-order and second-order classification of critical points in DLN loss landscapes. Classifies strict and non-strict saddles.
  \item \href{https://arxiv.org/abs/1910.01671}{Pure and Spurious Critical Points: a Geometric Study of Linear Networks} Matthew Trager et al. --- Defines spurious minima and counts the number of connected components of the global minima
  \item \href{https://arxiv.org/abs/2411.19920}{Geometry of fibers of the multiplication map of deep linear neural networks} Simon Pepin Lehalleur et al. --- Stratifies global minima into orbits using quiver representation theory
  \item \href{https://arxiv.org/abs/1412.0233}{The Loss Surfaces of Multilayer Networks} Anna Choromanska et al --- No-local-minima results in non-linear multilayer networks under some strong assumptions (spin glass models)
  \item \href{https://arxiv.org/abs/1802.10026}{Loss Surfaces, Mode Connectivity, and Fast Ensembling of DNNs}, Timur Garipov et al --- Study the connectedness of global minima of loss landscapes (mode connectivity)
  \item \href{https://arxiv.org/abs/1712.10132}{The Multilinear Structure of ReLU Networks}, Thomas Laurent, James von Brecht --- Show that local minima are typically singular
  \item Classic result: \href{https://www.sciencedirect.com/science/article/abs/pii/0893608089900142}{Neural networks and principal component analysis: Learning from examples without local minima} Pierre Baldi et al. --- Original result that linear network landscapes have no spurious local minima (single hidden layer case)
  \item Good review of key DLNs results: \href{https://arxiv.org/abs/2511.10362}{Gradient Flow Equations for Deep Linear Neural Networks: A Survey from a Network Perspective}, Joel Wendin, Claudio Altafini --- Also covers gradient flow
\end{itemize}

\subsection*{Implicit biases of gradient flow}
\begin{itemize}
  \item Saxe et al. \href{https://www.pnas.org/doi/10.1073/pnas.1820226116}{A mathematical theory of semantic development in deep neural networks} --- Read the supplementary materials section to understand the exact solution of gradient flow dynamics in deep linear networks
  \item \href{https://arxiv.org/abs/2106.15933}{Saddle-to-Saddle Dynamics in Deep Linear Networks: Small Initialization Training, Symmetry, and Sparsity} Arthur Jacot et al. --- Studies the saddle-to-saddle dynamics in DLNs. Introduces the different regimes (lazy and rich)
  \item \href{https://arxiv.org/abs/2411.09004}{The geometry of the deep linear network} Govind Menon --- Beautiful mathematical treatment of gradient flow in DLNs and a surprising mathematical result relating gradient flow and free-energy minimization. Makes implicit regularization explicit (under balanced assumptions)
  \item \href{https://arxiv.org/abs/2512.20607}{Saddle-to-Saddle Dynamics Explains A Simplicity Bias Across Neural Network Architectures} Saxe et al. --- Simplicity bias of gradient flow from DLNs extends to non-linear and transformer architectures under small initialization
  \item \href{https://arxiv.org/abs/2307.00144}{Abide by the Law and Follow the Flow: Conservation Laws for Gradient Flows} --- Gives a recipe to exhaustively derive conservation laws through gradient flow
  \item \href{https://arxiv.org/abs/2302.11055}{SGD learning on neural networks: leap complexity and saddle-to-saddle dynamics} --- SGD builds up complex solutions by first learning simpler solutions and then composing them
  \item \href{https://arxiv.org/abs/2405.17580}{Mixed Dynamics In Linear Networks: Unifying the Lazy and Active Regimes} --- Unifies lazy and rich in 2-layer DLNs and gives conditions for the transition between them
  \item \href{https://arxiv.org/abs/1806.07572}{Neural Tangent Kernel: Convergence and Generalization in Neural Networks} Arthur Jacot, Franck Gabriel, Clément Hongler --- NTK paper that describes gradient flow in function space
  \item \href{https://arxiv.org/abs/1812.07956}{On Lazy Training in Differentiable Programming} Lénaïc Chizat, Edouard Oyallon, Francis Bach --- Lazy training can occur in various settings, and is caused by scaling choices
  \item \href{https://arxiv.org/abs/2207.10430}{The Neural Race Reduction: Dynamics of Abstraction in Gated Networks} Saxe et al. --- Implicit biases toward shared representation in non-linear networks as modelled by gated DLNs
  \item \href{https://arxiv.org/abs/2506.06489}{Alternating Gradient Flows: A Theory of Feature Learning in Two-layer Neural Networks} Daniel Kunin et al. --- A theory of feature learning as a two-step process between dormant and active neurons
  \item \href{https://arxiv.org/abs/2409.14623}{From Lazy to Rich: Exact Learning Dynamics in Deep Linear Networks} --- Studies the transition between lazy and rich regimes in DLNs with balancedness parameters
  \item \href{https://arxiv.org/abs/2406.06158}{Get rich quick: exact solutions reveal how unbalanced initializations promote rapid feature learning} Daniel Kunin et al. --- (Lack of) balance between layers at initialization is preserved under GF, and influences rich vs lazy and implicit bias
  \item \href{https://arxiv.org/abs/1705.09280}{Implicit Regularization in Matrix Factorization} --- Gradient descent on matrix factorization converges to the minimum nuclear norm
  \item \href{https://arxiv.org/abs/1906.05890}{Gradient Descent Maximizes the Margin of Homogeneous Neural Networks} --- Implicit bias toward max-margin in classification
  \item \href{https://arxiv.org/abs/1810.02281}{A Convergence Analysis of Gradient Descent for Deep Linear Neural Networks} Sanjeev Arora et al. --- Convergence analysis of SGD
\end{itemize}

\subsection*{Learning rate (discrete GD)}
\begin{itemize}
  \item \href{https://arxiv.org/abs/2209.15594}{Self-Stabilization: The Implicit Bias of Gradient Descent at the Edge of Stability} --- Shows that curvature stabilizes around twice the inverse of the learning rate (2/eta).
  \item \href{https://arxiv.org/abs/2410.24206}{Understanding Optimization in Deep Learning with Central Flows} --- GD with a discrete learning rate is equivalent to gradient flow on an effective loss
  \item \href{https://arxiv.org/abs/2410.05192}{Understanding Warmup-Stable-Decay Learning Rates: A River Valley Loss Landscape Perspective} --- Pretraining exhibits a river valley loss landscape and gives intuition about the learning-rate schedule (warm-up, stable, decay)
  \item \href{https://www.nature.com/articles/s41467-025-58532-9}{Optimization on multifractal loss landscapes explains a diverse range of geometrical and dynamical properties of deep learning} --- Models the landscape as multifractal and analyses sub- and super-diffusive behaviour of GD
\end{itemize}

\subsection*{Stochasticity}
\begin{itemize}
  \item \href{https://arxiv.org/abs/2602.12257}{On the implicit regularization of Langevin dynamics with projected noise} --- Makes explicit the implicit bias of stochasticity in deep linear networks with a Langevin model
  \item \href{https://arxiv.org/abs/2306.08590}{Beyond Implicit Bias: The Insignificance of SGD Noise in Online Learning} --- Golden Path Hypothesis: in the transient regime dominated by drift (typical of pretraining), SGD is simply a noisy deformation of GD (it does not select different basins)
  \item \href{https://arxiv.org/abs/2306.04251}{Stochastic Collapse: How Gradient Noise Attracts SGD Dynamics Towards Simpler Subnetworks} --- Stochasticity induces a bias toward simpler solutions in deep linear networks
  \item \href{https://arxiv.org/abs/2109.14119}{Stochastic Training is Not Necessary for Generalization} --- Makes the implicit regularization of stochasticity explicit
  \item \href{https://arxiv.org/abs/2106.09524}{Implicit Bias of SGD for Diagonal Linear Networks: a Provable Benefit of Stochasticity} --- Shows that stochasticity has a bias toward sparser solutions relative to GD
  \item \href{https://arxiv.org/abs/1710.11029}{Stochastic gradient descent performs variational inference, converges to limit cycles for deep networks} --- SGD minimizes some free energy and can have circular currents at convergence
  \item \href{https://arxiv.org/abs/1704.04289}{Stochastic Gradient Descent as Approximate Bayesian Inference} --- Classic paper on SGD locally approximating Bayesian inference, although it assumes non-degeneracies
  \item \href{https://arxiv.org/abs/2002.03495}{A Diffusion Theory For Deep Learning Dynamics: Stochastic Gradient Descent Exponentially Favors Flat Minima} --- Implicit bias toward flat minima
  \item \href{https://arxiv.org/abs/2205.07069}{Implicit Regularization or Implicit Conditioning? Exact Risk Trajectories of SGD in High Dimensions} --- Shows the Golden Path Hypothesis on quadratic loss in a convex setup using an interesting SDE model of SGD.
  \item \href{https://arxiv.org/abs/1812.06162}{An Empirical Model of Large-Batch Training} --- Gradient noise scale can be used to decide the optimal batch size
  \item \href{https://arxiv.org/abs/2503.22478}{Almost Bayesian: The Fractal Dynamics of Stochastic Gradient Descent} --- Relates SGD to Bayesian training
  \item \href{https://arxiv.org/abs/2306.04815}{Catapults in SGD: spikes in the training loss and their impact on generalization through feature learning} --- Loss spikes when training with SGD impact generalization
  \item \href{https://arxiv.org/abs/2006.04740}{The Heavy-Tail Phenomenon in SGD} --- SGD noise can be heavy-tailed, and this induces a different regularization
\end{itemize}

\subsection*{Emergence: empirical examples}
\begin{itemize}
  \item \href{https://arxiv.org/abs/2201.02177}{Grokking: Generalization Beyond Overfitting on Small Algorithmic Datasets} --- Grokking: delayed generalization on modular addition
  \item \href{https://transformer-circuits.pub/2022/in-context-learning-and-induction-heads/index.html}{In-context Learning and Induction Heads} --- Induction heads form during training. They are important for in-context learning
  \item \href{https://arxiv.org/abs/2502.17424}{Emergent Misalignment: Narrow finetuning can produce broadly misaligned LLMs} --- Emergent misalignment: fine-tuning on insecure code can generalize to misaligned behaviour on other data
  \item \href{https://arxiv.org/abs/2602.07852}{Emergent Misalignment is Easy, Narrow Misalignment is Hard} --- By regularizing, we can mitigate emergent misalignment
  \item \href{https://arxiv.org/abs/2111.00034}{Neural Networks as Kernel Learners: The Silent Alignment Effect} --- While the loss is flat, student singular vectors align with teacher singular vectors, and this can be detected with the NTK
  \item \href{https://arxiv.org/abs/2304.15004}{Are Emergent Abilities of Large Language Models a Mirage?} --- Emergence is in the eye of the metric used to measure it
  \item \href{https://arxiv.org/abs/2203.15556}{Training Compute-Optimal Large Language Models} --- Chinchilla scaling law paper
  \item \href{https://arxiv.org/abs/2206.07682}{Emergent Abilities of Large Language Models} --- Shows that novel capabilities emerge with more compute
\end{itemize}

\subsection*{Theoretical approaches to emergence}
\begin{itemize}
  \item \href{https://arxiv.org/abs/2310.06110}{Grokking as the Transition from Lazy to Rich Training Dynamics}
  \item \href{https://arxiv.org/abs/2310.03789}{Grokking as a First Order Phase Transition in Two Layer Networks, Rubin et al.}
  \item \href{https://www.youtube.com/watch?v=WcWFFiPRslM&t=1056s}{Blake Bordelon - Infinite limits and scaling laws of neural networks - IPAM at UCLA}
  \item \href{https://arxiv.org/abs/2307.15936}{A Theory for Emergence of Complex Skills in Language Models}
  \item \href{https://ericjmichaud.com/quanta/}{On neural scaling and the quanta hypothesis}
  \item \href{https://pehlevan.seas.harvard.edu/sites/g/files/omnuum6471/files/pehlevan/files/princeton_lecture_notes.pdf}{Lecture Notes on Infinite-Width Limits of Neural Networks; Cengiz Pehlevan and Blake Bordelon}
  \item \href{https://arxiv.org/abs/2601.01010}{Disordered Dynamics in High Dimensions: Connections to Random Matrices and Machine Learning; Blake Bordelon, Cengiz Pehlevan}
  \item \href{https://arxiv.org/abs/2502.18553}{Applications of Statistical Field Theory in Deep Learning; Zohar Ringel et al.}
  \item \href{https://link.springer.com/book/10.1007/978-3-030-46444-8}{Statistical Field Theory for Neural Networks, Moritz Helias, David Dahmen}
  \item \href{https://arxiv.org/abs/2602.12855}{Lecture notes: From Gaussian processes to feature learning, Moritz Helias et al}
\end{itemize}


\bibliographystyle{alphaurl}
\bibliography{biblo}

\end{document}
